多模态大语言模型（MLLM）在跨模态理解与生成上已取得显著进展，但仍持续受幻觉问题困扰，即生成与视觉输入不一致的文本。我们指出，其根本原因之一在于自回归解码阶段的\textbf{跨模态粒度失配}：图像被编码为连续的整体特征，而文本必须以离散、且常呈语义碎片化的子词（sub-word）token 序列输出；这使得难以将无独立语义的单 token 可靠地锚定到具体视觉实体，模型易进入未经验证的解码状态，语言先验的统计惯性进而放大幻觉。现有缓解策略往往依赖昂贵重训或高代价的多轮验证。为此，我们提出 \textbf{CHORD}（Calibrating Hallucinations via Object-Resonant Decoding，通过物体共振解码校准幻觉）——一种\textbf{免训练}、可即插即用的解码框架，借助 \textbf{Object-Resonant Decoding（物体共振解码）} 将生成轨迹重新锚定到视觉事实。CHORD 在活跃解码前沿上工作，动态评估\textbf{候选文本片段}与\textbf{外部视觉证据}的对齐程度：仅当片段与具体视觉物体在语义上``共振''时增强其概率，而对仅由语言惯性驱动的无根据内容予以抑制。为在控制开销的同时捕获\textbf{过度自信型}幻觉，我们采用\textbf{连续 Top-$1$ 监控}并在必要时\textbf{回退到全 $K$ 候选}的共振计算，并利用外部文本编码器的 KV 缓存复用，使完整前沿评分仅在 Top-$1$ 假设未通过共振检验时触发。在多个面向幻觉的基准上，CHORD 显著降低物体幻觉；作为与宿主架构解耦的框架，它无需额外训练，为可信 MLLM 部署提供可扩展路径。
